\documentclass[aspectratio=169, 10pt]{beamer}
\usetheme{Madrid}
\usecolortheme{default}

\usepackage[utf8]{inputenc}
\usepackage{amsmath, amssymb, amsthm}
\usepackage{bm}
\usepackage{graphicx}
\usepackage{xcolor}

\title{Uncertainty Estimates of Predictions via a General Bias-Variance Decomposition}
\author{Sebastian G. Gruber, Florian Buettner}
\date{AISTATS 2023}

\begin{document}

% ---------- SLIDE 1: TITLE ----------
\begin{frame}
    \titlepage
    \begin{center}
        \tiny{Presented by: Vadim Kasiuk}
    \end{center}
\end{frame}

% ---------- SLIDE 2: THE PROBLEM & MOTIVATION ----------
\begin{frame}{The Problem: Uncertainty Estimation Under Domain Drift}
    \begin{itemize}
        \item \textbf{Context:} In safety-critical applications (e.g., cancer diagnostics), knowing \emph{when} a model is wrong is as important as being correct.
        \item \textbf{Standard Practice:} Use the maximum softmax probability (confidence score). However, this is unreliable under \textbf{domain drift} (Ovadia et al., 2019).
        \item \textbf{Limitation of Existing Theory:} Proper scores (like Log-Likelihood) capture uncertainty, but lack a general bias-variance decomposition to quantify the \emph{variance} of the prediction.
    \end{itemize}
    \vfill
    \begin{block}{Research Question}
        Can we derive a general bias-variance decomposition for strictly proper scoring rules where the variance term naturally captures predictive uncertainty?
    \end{block}
\end{frame}

% ---------- SLIDE 3: RELATED WORK & GAP ----------
\begin{frame}{Related Work \& Identified Gaps}
    \begin{table}[]
        \centering
        \begin{tabular}{p{0.4\linewidth} p{0.5\linewidth}}
            \hline
            \textbf{Method / Work} & \textbf{Limitation / Gap} \\
            \hline
            James \& Hastie (1997) & General decomposition, but \textbf{no closed-form solution}. \\
            Heskes (1998) / Hansen \& Heskes (2000) & Specific to KL-divergence \& Exponential Families only. \\
            Domingos (2000) & Unclear if decomposition holds for arbitrary losses. \\
            \textbf{This work (Gruber \& Buettner)} & $\checkmark$ General closed-form for proper scores. \\
            & $\checkmark$ Unifies previous decompositions. \\
            & $\checkmark$ Works in \textbf{logit space} (no normalization required). \\
            \hline
        \end{tabular}
    \end{table}
\end{frame}

% ---------- SLIDE 4: BACKGROUND - BREGMAN DIVERGENCES ----------
\begin{frame}{Background: Bregman Divergences}
    \begin{definition}[Bregman Divergence]
        Given a differentiable, convex function $\phi: U \to \mathbb{R}$, the Bregman divergence is:
        \[
        d_\phi(x, y) = \phi(y) - \phi(x) - \langle \nabla \phi(x), y - x \rangle
        \]
        Interpretation: Difference between $\phi(y)$ and the tangent at $\phi(x)$.
    \end{definition}
    \begin{itemize}
        \item Non-negative: $d_\phi \ge 0$, equality iff $x = y$.
        \item Generates Squared Error: If $\phi(x) = x^2$, then $d_\phi(x,y) = (x-y)^2$.
    \end{itemize}
    
\end{frame}

% ---------- SLIDE 5: BACKGROUND - PROPER SCORES ----------
\begin{frame}{Background: Strictly Proper Scoring Rules}
    \begin{definition}[Strictly Proper Score]
        A scoring rule $S(P, y)$ is \emph{strictly proper} if the expected score is maximized iff the predicted distribution $P$ equals the true data distribution $Q$:
        \[
        \mathbb{E}_{y \sim Q}[S(Q, y)] \ge \mathbb{E}_{y \sim Q}[S(P, y)] \quad \forall P
        \]
    \end{definition}
    \begin{itemize}
        \item \textbf{Examples:} Log-Likelihood ($S(P,y)=\log p(y)$), Brier Score, Continuous Ranked Probability Score (CRPS).
        \item \textbf{Relation to Bregman:} Strictly proper scores correspond to Bregman divergences (Ovcharov, 2018).
        \[
        G(Q) - S(P) \cdot Q = d_{G,S}(P, Q)
        \]
        where $G$ is the negative entropy.
    \end{itemize}
\end{frame}

% ---------- SLIDE 6: CORE CONTRIBUTION - BIAS-VARIANCE DECOMPOSITION ----------
\begin{frame}{Main Result: General Bias-Variance Decomposition}
    \begin{theorem}[General BVD for Proper Scores]
        For a strictly proper score $S$ with negative entropy $G$, estimated prediction $\hat{f}$, and target $Y \sim Q$:
        \[
        \mathbb{E}[-S(\hat{f})(Y)] = \underbrace{-G(Q)}_{\text{Noise}} + \underbrace{B_{G^*}[S(\hat{f})]}_{\text{Variance}} + \underbrace{d_{G^*, S^{-1}}\big(S(Q), \mathbb{E}[S(\hat{f})]\big)}_{\text{Bias}^2}
        \]
    \end{theorem}
    \begin{itemize}
        \item $B_{G^*}$ is the \textbf{Bregman Information} generated by the convex conjugate $G^*$.
        \item The variance term generalizes the classical variance $\mathbb{V}[X]$ (since $B_{x^2}[X] = \mathbb{V}[X]$).
    \end{itemize}

    \begin{definition}[Bregman Information (Banerjee et al., 2005)]
        For a convex function $\phi$ and random variable $X$:
        \[
        B_\phi[X] = \mathbb{E}[d_\phi(\mathbb{E}[X], X)] = \mathbb{E}[\phi(X)] - \phi(\mathbb{E}[X])
        \]
        \textit{Interpretation:} Measures the \textbf{spread} of $X$ around its mean. Generalization of Variance.
    \end{definition}
    
\end{frame}

% ---------- SLIDE 7: BREGMAN INFORMATION (THE "VARIANCE") ----------
\begin{frame}{The Variance Term: Bregman Information}
    
    \begin{figure}[h]
        \centering
        \includegraphics[width=0.5\linewidth]{breg_inf.jpeg}
    \end{figure}
\end{frame}

% ---------- SLIDE 8: SPECIAL CASE - EXPONENTIAL FAMILIES ----------
\begin{frame}{Special Case: Exponential Families}
    \begin{block}{Proposition (Exponential Family)}
        For an exponential family $p_\theta(y) = \exp(\langle \theta, T(y) \rangle - A(\theta)) h(y)$:
        \[
        \mathbb{E}[-\ln p_{\hat{\theta}}(Y)] = n(\theta) + B_A[\hat{\theta}] + d_A(\theta, \mathbb{E}[\hat{\theta}])
        \]
        where $A$ is the log-partition function.
    \end{block}
    \textbf{Implication:}
    \begin{itemize}
        \item The variance term is $B_A[\hat{\theta}]$ (classical variance in the \textbf{natural parameter space} $\theta$).
        \item Recovers the standard MSE decomposition (Example 3.5).
    \end{itemize}
\end{frame}

% ---------- SLIDE 9: SPECIAL CASE - CLASSIFICATION (LOGIT SPACE) ----------
\begin{frame}{Special Case: Classification (Logit Space)}
    \begin{corollary}[Classification NLL Decomposition]
        For logits $\hat{z} \in \mathbb{R}^k$ and softmax $\text{sm}$:
        \[
        \mathbb{E}[-\ln \text{sm}_Y(\hat{z})] = H(Q) + B_{\text{LSE}}[\hat{z}] + d_{\text{LSE}}(\text{sm}^{-1}(Q), \mathbb{E}[\hat{z}])
        \]
        where $\text{LSE}(x) = \ln \sum e^{x_i}$ is the LogSumExp function.
    \end{corollary}
    \begin{itemize}
        \item \textbf{Surprising result:} The variance $B_{\text{LSE}}[\hat{z}]$ measures uncertainty directly in the \textbf{logit space}.
        \item No normalization to probabilities required! This is highly practical for Deep Learning.
    \end{itemize}
\end{frame}

% ---------- SLIDE 10: APPLICATIONS - WHY ENSEMBLES WORK ----------
\begin{frame}{Application 1: Why Ensembles Reduce Uncertainty}
    \begin{itemize}
        \item \textbf{Law of Total Variance for Bregman Information:}
        \[
        B_G[X] = \mathbb{E}[B_G[X \mid Y]] + B_G[\mathbb{E}[X \mid Y]]
        \]
        \item \textbf{Interpretation for Ensembles:} Ensemble prediction $\hat{\theta}_D^{(n)}$ marginalizes out source of randomness $W$ (e.g., initialization).
        \item As $n \to \infty$, $B_A[\hat{\theta}_D^{(n)}] \to B_A[\mathbb{E}[\theta_{W,D} \mid D]]$ (the variance due to $W$ vanishes).
        \item \textbf{Conclusion:} Ensembles strictly improve expected performance (Equation 5).
    \end{itemize}

    \begin{itemize}
        \item Using Markov's inequality on the Bregman divergence:
        \[
        P\left(d_G(\mathbb{E}[X], X) \ge \frac{1}{\alpha} B_G[X]\right) \le \alpha
        \]
        \item \textbf{Confidence Region:}
        \[
        \text{CR}_{(1-\alpha)} = \left\{ x \mid d_G(\mathbb{E}[X], x) \le \frac{B_G[X]}{\alpha} \right\}
        \]
    \end{itemize}
\end{frame}

% ---------- SLIDE 11: APPLICATIONS - CONFIDENCE REGIONS ----------
\begin{frame}{Application 2: Confidence Regions via Markov's Inequality}
    \begin{figure}[h]
        \centering
        \includegraphics[width=0.45\linewidth]{confid.jpeg}
        \includegraphics[width=0.45\linewidth]{iris_dist.jpeg}
    \end{figure}
\end{frame}

% ---------- SLIDE 12: EXPERIMENTS - CIFAR-10C / IMAGENET-C ----------
\begin{frame}{Experiments: Out-of-Distribution Detection}
    \begin{figure}[h]
        \centering
        \includegraphics[width=0.7\linewidth]{cifar-10.jpeg}
    \end{figure}
    \textbf{Key Finding:}
    \begin{itemize}
        \item To maintain validation accuracy under severe corruption (CIFAR-10C), we discard only $\sim 7\%$ of data using Bregman Info vs $\sim 14\%$ using confidence.
        \item Bregman Information is a superior measure of \textbf{instance-wise uncertainty} under domain drift.
    \end{itemize}
\end{frame}

% ---------- SLIDE 13: COMPARISON WITH BASELINES ----------
\begin{frame}{Results: Common Classifiers on Toy Data}
    \begin{itemize}
        \item \textbf{Setup:} Sampled multiple training sets to approximate ground truth Bregman Information (Section 4.1).
        \item \textbf{Observation:}
        \begin{itemize}
            \item \textit{High Bregman Info:} Decision Trees, k-NN (high variance due to data sensitivity).
            \item \textit{Low Bregman Info:} Random Forests, XGBoost (ensemble methods reduce variance).
            \item Neural Networks exhibit moderate Bregman Info but are sensitive to initialization.
        \end{itemize}
        \item \textbf{Qualitative insight:} The decomposition aligns with the intuitive notion of "model stability."
    \end{itemize}
\end{frame}

% ---------- SLIDE 14: DISCUSSION & LIMITATIONS ----------
\begin{frame}{Discussion \& Limitations}
    \begin{block}{Strengths}
        \begin{itemize}
            \item Unifies bias-variance theory under a single mathematical framework (Bregman).
            \item Provides practical tools (ensembles, confidence regions, OOD detection).
            \item Logit-space formulation is ideal for modern neural networks.
        \end{itemize}
    \end{block}
    \begin{block}{Limitations \& Future Work}
        \begin{itemize}
            \item \textbf{Computation:} Estimating Bregman Information requires multiple forward passes or an ensemble (computationally expensive).
            \item \textbf{Assumption:} Strictly proper scores only. Does not apply to 0-1 loss.
            \item \textbf{Future:} Extending to Bayesian Neural Networks and uncertainty quantification in LLMs.
        \end{itemize}
    \end{block}
\end{frame}

% ---------- SLIDE 15: CONCLUSION ----------
\begin{frame}{Conclusion}
    \begin{itemize}
        \item We introduced the first general \textbf{bias-variance decomposition for strictly proper scores}.
        \item The \textbf{Bregman Information} emerges as the universal variance term, generalizing classical variance.
        \item We derived practical formulations for \textbf{exponential families} and, crucially, \textbf{classification in the logit space} (LogSumExp).
        \item Demonstrated utility in:
        \begin{itemize}
            \item Explaining the success of \textbf{ensembles}.
            \item Building \textbf{confidence regions}.
            \item Reliable \textbf{out-of-distribution detection} under domain drift.
        \end{itemize}
    \end{itemize}
    \vfill
    \centering
    \textbf{Code:} \url{https://github.com/MLO-lab/Uncertainty_Estimates_via_BVD}
\end{frame}

\end{document}
